\subsection*{Cap\'itulo 1: Fundamentos matem\'aticos}
\textbf{1.1}\label{resp:1-1}\quad
(1) $A\cup B=\{1, 2, 3, 4, 5\}$, $A\cap B=\{3, 4\}$, $A\setminus B=\{1, 2\}$. (2) $\{x\in\mathbb Z^+\mid x\text{ par}\wedge x\le10\}$. (3) $2^5=32$. (4) $A^c=\{6, 7\}$; $A\,\Delta\,B=\{1, 2, 3, 6\}$. (5) $(A\cup B)^c=\{4\}$ y $A^c\cap B^c=\{3, 4\}\cap\{1, 4\}=\{4\}$. (6) $30+25-10=45$. (7) $120+90-40=170$ usan alguna; no usan ninguna: $30$. (8) El primero s\'i; el segundo no (no son disjuntos: comparten el 2).

\medskip\noindent\textbf{1.2}\label{resp:1-2}\quad
(1) No, s\'i (V), no (predicado abierto). (2) V, F, V, V. (3) ``Existe un usuario sin token v\'alido''. (4) Ambas columnas dan V, F, V, V. (5) $\neg p\vee q$. (6) Tautolog\'ia, contradicci\'on, contingencia. (7) Modus Tollens con $q\to r$ y $\neg r$ da $\neg q$; Modus Tollens con $p\to q$ y $\neg q$ da $\neg p$. (8) \texttt{if (!a || b)}.

\medskip\noindent\textbf{1.3}\label{resp:1-3}\quad
(1) 6. (2) S\'i. (3) No (falta $(2, 1)$). (4) $\begin{bmatrix}0&1&0\\0&0&1\\0&0&1\end{bmatrix}$. (5) Se agregan $(1, 3),(2, 4),(1, 4)$. (6) S\'i: en el Hasse, 1 abajo, 2 y 3 en medio, 6 arriba. (7) Reflexiva, sim\'etrica y transitiva; 3 clases ($\{0, 3, 6\}$, $\{1, 4, 7\}$, $\{2, 5, 8\}$). (8) Todo el cierre son unos: $\begin{bmatrix}1&1&1\\1&1&1\\1&1&1\end{bmatrix}$ (el ciclo $1\to2\to3\to1$).

\medskip\noindent\textbf{1.4}\label{resp:1-4}\quad
(1) 7. (2) 3, 4 y 2. (3) No (porque $f(2)=f(-2)$); en $[0,\infty)$ s\'i. (4) 9 y 25. (5) $f^{-1}(x)=\frac{x+6}{3}$. (6) $5n+3\le5n+3n=8n$ para $n\ge1$: $c=8$, $n_0=1$. (7) $\log n,\ n,\ n\log n,\ n^2,\ 2^n$. (8) $4n^2+10n+7\le21n^2$ para $n\ge1$: $c=21$, $n_0=1$.

\medskip\noindent\textbf{1.5}\label{resp:1-5}\quad
(1) 7. (2) $6xy$. (3) $(2, 4)$. (4) $x^2+y^2\le25$ (c\'irculo de radio 5). (5) $3x^2y^2-4$ y $2x^3y+5$. (6) Circunferencia de radio 2 centrada en el origen. (7) $(3{,}2;\ 1{,}6)$. (8) $0{,}5c+30>50\Rightarrow c>40$.

\medskip\noindent\textbf{Mini proyecto:} (1) $P(u)\equiv(u\in A)\wedge(V(u)=1)\wedge(R(c, i)\le50)$. (2) El bloqueo es $\neg P(u)\equiv(u\notin A)\vee(V(u)=0)\vee(R(c, i)>50)$; el bloqueo autom\'atico del servidor es simplemente $R(c, i)>50$. (3) $\partial R/\partial c=0{,}5$ y $\partial R/\partial i=10$: un intento fallido pesa 20 veces m\'as que un punto de CPU.

\medskip\noindent\textbf{Autoevaluaci\'on 1:}\label{resp:q1}\quad 1 b, 2 a, 3 b, 4 b, 5 c, 6 b, 7 c, 8 b, 9 b, 10 b.
